% !TEX root = ../main.tex

\chapter{实验求解流程}\label{ch:algorithm}
  第\ref{ch:methodology}章已经给出了线性方程与半线性方程对应的离散化模型和优化问题.
  为了便于后续实验分析, 本章从算法执行流程的角度总结本文方法, 重点说明样本路径生成,
  随机特征构造, 参数求解与测试评估之间的关系.

  \section{整体思路}\label{sec:overview}
    本文的实验流程可以概括为四个阶段:
    \begin{enumerate}
      \item 根据给定方程生成离散样本路径;
      \item 在所有时间层和路径点上构造随机特征, 用于近似 $Z_{k}$;
      \item 根据方程类型选择线性最小二乘或非线性 Levenberg--Marquardt 方法求解参数;
      \item 在独立测试路径上重新推进终端值, 计算测试残差.
    \end{enumerate}

    上述四步中, 前两步构成公共的数据准备阶段, 对线性和半线性问题都相同;
    第三步则是两类方程的主要差异所在.
    在线性情形下, 末端匹配可以直接化为一个最小二乘系统;
    在半线性情形下, 则需要通过迭代方式不断修正参数.
    算法的流程图如图~\ref{fig:flowchart} 给出.
    对应的较详细伪代码见附录~\ref{ch:pseudocode}.

    \begin{figure}[ht]
      \centering
      \includegraphics[width=0.92\textwidth]{flowchart}
      \caption{实验求解方法的整体流程}
      \label{fig:flowchart}
    \end{figure}

  \section{样本路径与随机特征的构造}\label{sec:sampling}
    设终值问题已经被离散为 $N$ 个时间步, 样本数为 $S$.
    对每个实验算例, 首先按照对应随机微分方程生成 $S$ 条独立样本路径
    \begin{equation*}
      \{X_{k}^{(i)}\}_{k = 0}^{N}, \qquad i = 1,2,\dots,S,
    \end{equation*}
    以及与之对应的布朗增量
    \begin{equation*}
      \{\Delta W_{k}^{(i)}\}_{k = 0}^{N - 1}.
    \end{equation*}
    这些路径同时承担两个作用: 一方面它们给出离散 BSDE 的推进轨道,
    另一方面它们也是随机特征映射的输入样本.

    在得到路径后, 对每个时间层和路径点计算随机特征
    \begin{equation*}
      H_{k}^{(i)} = \tanh(A x_{k}^{(i)} + b t_{k} + c),
    \end{equation*}
    其中 $A, b, c$ 为固定随机参数.
    本文始终只训练输出层参数 $\alpha$, 并用
    \begin{equation*}
      Z_{k}^{(i)} \approx \alpha H_{k}^{(i)}
    \end{equation*}
    作为梯度项的统一近似.
    于是原问题的求解就转化为对 $y_{0}$ 和 $\alpha$ 的求解.

  \section{线性方程的求解流程}\label{sec:exp_linear_flow}
    对线性方程, 由第\ref{ch:methodology}章的推导可知,
    每条路径的终端值都可以写成
    \begin{equation*}
      y_{N}^{(i)} = \beta_{i} y_{0} + \langle \Gamma_{i}, \alpha \rangle + \eta_{i}.
    \end{equation*}
    因而所有路径拼接后, 就得到一个关于参数向量 $(y_{0}, \operatorname{vec}(\alpha))$ 的线性最小二乘问题.

    从实验流程的角度看, 线性情形可分为以下几步:
    \begin{enumerate}
      \item 生成样本路径并计算随机特征;
      \item 根据离散递推式计算每条路径的系数 $\beta_{i}$, $\Gamma_{i}$, $\eta_{i}$;
      \item 将所有路径系数拼接成统一的设计矩阵;
      \item 求解最小二乘问题得到 $y_{0}$ 与 $\alpha$;
      \item 用求得的参数回代, 计算训练样本上的终端残差.
    \end{enumerate}

    线性流程的特点在于参数只需求解一次, 不需要额外的外层迭代.
    因此实验中线性部分主要用于观察样本数, 特征宽度以及维度变化对最小二乘近似效果的影响.

  \section{半线性方程的求解流程}\label{sec:exp_semilinear_flow}
    对半线性方程, 终端值关于参数不再保持线性, 因而不能直接写成单个最小二乘系统.
    这时需要把参数记为
    \begin{equation*}
      \theta = \bigl(y_{0}, \operatorname{vec}(\alpha)\bigr),
    \end{equation*}
    并通过最小化终端残差平方和来求解
    \begin{equation*}
      \min_{\theta} \frac{1}{2}\lVert r(\theta) \rVert_{2}^{2}.
    \end{equation*}

    本文实验中对半线性方程采用 Levenberg--Marquardt 方法.
    每次迭代都围绕当前参数执行以下流程:
    \begin{enumerate}
      \item 根据当前参数构造 $Z_{k}^{(i)} = \alpha H_{k}^{(i)}$;
      \item 沿所有样本路径前向推进离散 BSDE, 得到终端残差;
      \item 计算残差关于参数的雅可比矩阵;
      \item 求解阻尼线性子问题, 得到参数增量;
      \item 比较更新前后的目标函数值, 决定是否接受该步并调整阻尼参数.
    \end{enumerate}

    这一步骤会持续进行, 直到终端残差足够小, 参数更新足够小, 或达到最大迭代次数.
    因而半线性实验的核心并不在于单次前向传播本身,
    而在于“前向残差评估 + 雅可比构造 + 阻尼更新”这一循环是否稳定.

  \section{测试阶段与误差评估}\label{sec:test&eval}
    为了避免只在固定训练路径上评价模型, 本文在训练完成后还会重新采样独立测试路径.
    对测试路径仍然使用相同的随机特征参数和输出层参数,
    并按照与训练阶段一致的离散格式推进终端值.
    由此得到的测试误差记为
    \begin{equation*}
      \operatorname{RMSE}_{\text{test}}
      =
      \sqrt{
        \frac{1}{S}
        \sum_{i = 1}^{S}
        \bigl(
          Y_{N}^{(i)} - g(X_{N}^{(i)})
        \bigr)^{2}
      }.
    \end{equation*}

    若训练误差下降而测试误差没有同步下降, 则说明随机特征近似可能更偏向于拟合当前训练路径;
    反之, 若两者同时下降, 则说明该参数设置具有更稳定的泛化效果.
    因而训练 RMSE 与测试 RMSE 的对比是本实验章节中的核心观察指标之一.
